\documentclass[12pt, a4paper, oneside]{ctexart}
\usepackage{amsmath, amsthm, amssymb, bm, color, framed, graphicx, hyperref, mathrsfs}

\title{\textbf{3月28日作业}}
\author{韩岳成 524531910029}
\date{\today}
\linespread{1.5}
\definecolor{shadecolor}{RGB}{241, 241, 255}
\newcounter{problemname}
\newenvironment{problem}{\begin{shaded}\stepcounter{problemname}\par\noindent\textbf{题目\arabic{problemname}. }}{\end{shaded}\par}
\newenvironment{solution}{\par\noindent\textbf{解答. }}{\par}
\newenvironment{note}{\par\noindent\textbf{题目\arabic{problemname}的注记. }}{\par}
\everymath{\displaystyle}

\begin{document}

\maketitle

\begin{problem}
    若一个电子的动能等于静能，试求该电子的速率和德布罗意波长。
\end{problem}

\begin{solution}
    电子动能：
    \[E_k = (\gamma - 1) m c^2\]
    电子静能：
    \[E_0 = m c^2\]
    由题意可得：
    \[(\gamma - 1) m c^2 = m c^2 \Rightarrow \gamma = 2\]
    由此可得电子的速率：
    \[v = c \sqrt{1 - \frac{1}{\gamma^2}} = \frac{\sqrt{3}}{2} c\]
    由$E^2 = (pc)^2 + (m c^2)^2$与 $E = mc^2 + E_k = 2mc^2$可得电子的动量：
    \[p = \sqrt{3}m c\]
    故德布罗意波长
    \[ \lambda = \frac{h}{p} = \frac{h}{\sqrt{3} m c}\]
\end{solution}

\begin{problem}
利用不确定关系估算氢原子基态能量及半径。
\end{problem}

\begin{solution}
    由不确定度关系$\Delta x \Delta p \sim \hbar$，设半径为$r$，则$\Delta x \sim r$，$\Delta p \sim \frac{\hbar}{r}$。
    
    氢原子基态能量为动能与势能之和：
    \[E = \frac{(\Delta p)^2}{2m} - \frac{e^2}{4\pi \epsilon_0 r}\]
    将$\Delta p$代入上式：
    \[ E = \frac{\hbar^2}{2 m r^2} - \frac{e^2}{4\pi \epsilon_0 r} \]
    对$E$关于$r$求导并令其为零以找到能量的极小值：
    \[ \frac{dE}{dr} = -\frac{\hbar^2}{m r^3} + \frac{e^2}{4\pi \epsilon_0 r^2} = 0 \]
    解出$r$：
    \[ r = \frac{4\pi \epsilon_0 \hbar^2}{m e^2} \]
    将$r$代入能量表达式中：
    \[ E = -\frac{m e^4}{32 \pi^2 \epsilon_0^2 \hbar^2} \]
    这就是氢原子基态的能量和半径的估算值。 
\end{solution}

\begin{problem}
已知粒子波函数 $\Psi = N \exp\left\{-\frac{|x|}{2a} - \frac{|y|}{2b} - \frac{|z|}{2c}\right\}$，试求：
\begin{enumerate}
\item 归一化常数,
\item 粒子的 $x$ 坐标在 $0$ 到 $a$ 之间的概率,
\item 粒子的 $y$ 坐标和 $z$ 坐标分别在 $-b$ 到 $+b$ 和 $-c$ 到 $+c$之间的概率。
\end{enumerate}
\end{problem}

\begin{solution}
1. 
    归一化条件要求：
    \[ \int_{-\infty}^{\infty} \int_{-\infty}^{\infty} \int_{-\infty}^{\infty} |\Psi|^2 dx dy dz = 1 \]
    由于$\int_{-\infty}^{\infty} e^{-\frac{|x|}{a}} dx = 2a$，同理对于$y$和$z$坐标也有类似的积分结果，因此：
    \[ N^2 (2a)(2b)(2c) = 1 \Rightarrow N = \frac{1}{2\sqrt{2abc}} \]

2.
    粒子在$x$坐标在$0$到$a$之间的概率为：
    \[ P_x = \int_0^a |\Psi|^2 dx\]
    此时$x$的波函数部分是$\frac{1}{2\sqrt{a}}e^{-\frac{|x|}{a}}$，那么概率密度为$\frac{1}{2a}e^{-\frac{|x|}{a}}$，因此：
    \[ P(0<x<a) = \int_0^a \frac{1}{2a} e^{-\frac{x}{a}} dx = \frac{1}{2}(1 - e^{-1}) \]

3.
    由对称性，$y$坐标在$-b$到$+b$之间的概率为：
    \[ P_y = 1 - e^{-1} \]
    同样的，$z$坐标在$-c$到$+c$之间的概率也为：
    \[ P_z = 1 - e^{-1} \]
    由于$y$和$z$坐标的概率是独立的，因此它们同时满足条件的概率为：
    \[ P(-b<y<b, -c<z<c) = P_y \cdot P_z = (1 - e^{-1})^2 \]
\end{solution}

\begin{problem}
设质量为 $m$ 的粒子在半壁无限高的一维方势阱中运动，试求在 $E < V_0$ 的束缚态情况下：
\begin{enumerate}
    \item 粒子能级的表达式；
    \item 势阱深 $V_0$ 和势阱宽 $a$ 需满足什么关系，才能保证此阱中至少存在一个束缚态？
\end{enumerate}
\[ V = \begin{cases} \infty & x < 0 \\ 0 & 0 \le x \le a \\ V_0 & x > a \end{cases} \]
\end{problem}

\begin{solution}
1.
    在$0 \le x \le a$区域，粒子满足的薛定谔方程为：
    \[ -\frac{\hbar^2}{2m} \frac{d^2 \psi}{dx^2} = E \psi \]
    其解为：
    \[ \psi(x) = A \sin(kx) + B \cos(kx) \]
    其中$k = \sqrt{\frac{2mE}{\hbar^2}}$。由于在$x=0$处势能为无穷大，波函数必须满足$\psi(0) = 0$，因此$B=0$，波函数简化为：
    \[ \psi(x) = A \sin(kx) \]
    
    在$x > a$区域，粒子满足的薛定谔方程为：
    \[ -\frac{\hbar^2}{2m} \frac{d^2 \psi}{dx^2} + V_0 \psi = E \psi \]
    其解为：
    \[ \psi(x) = C e^{-\kappa (x-a)} + D e^{\kappa (x-a)}\]
    其中$\kappa = \sqrt{\frac{2m(V_0 - E)}{\hbar^2}}$。
    由于束缚态波函数必须在$x\to \infty$处衰减，因此$D=0$，波函数简化为：
    \[ \psi(x) = C e^{-\kappa (x-a)} \]

    在$x=a$处，波函数及其导数必须连续，因此有以下边界条件：
    \[ A \sin(ka) = C e^{-\kappa (a-a)} = C \]
    \[ k A \cos(ka) = -\kappa C e^{-\kappa (a-a)} = -\kappa C \]
    
    将第一个边界条件代入第二个边界条件中，得到能级的表达式：
    \[ k \cot(ka) = -\kappa \]
    这就是束缚态能级的表达式。

2.
由第 1 问可知，束缚态（$0 < E < V_0$）的波函数需满足超越方程：
\[
    k \cot(ka) = -\kappa
\]
其中，由于 $E > 0$ 且 $E < V_0$，波数 $k$ 和衰减常数 $\kappa$ 均为正实数：
\[
    k = \frac{\sqrt{2mE}}{\hbar} > 0, \quad \kappa = \frac{\sqrt{2m(V_0 - E)}}{\hbar} > 0
\]

为了方便讨论，我们在方程两边同乘势阱宽度 $a$，引入无量纲变量 $\xi$ 和 $\eta$：
\[
    \xi = ka > 0, \quad \eta = \kappa a > 0
\]
此时，超越方程化为：
\[
    \xi \cot(\xi) = -\eta \quad \Rightarrow \quad \eta = -\xi \cot(\xi)
\]

同时，注意 $\xi$ 和 $\eta$ 满足一个圆的方程约束。计算它们的平方和：
\[
    \xi^2 + \eta^2 = (ka)^2 + (\kappa a)^2 = \frac{2mEa^2}{\hbar^2} + \frac{2m(V_0 - E)a^2}{\hbar^2} = \frac{2mV_0a^2}{\hbar^2}
\]
令 $R = \sqrt{\frac{2mV_0a^2}{\hbar^2}}$，则 $\xi$ 和 $\eta$ 必须同时满足第一象限内的圆方程：
\[
    \xi^2 + \eta^2 = R^2 \quad (\text{其中 } \xi > 0, \eta > 0)
\]

问题的本质转化为：\textbf{在 $\xi-\eta$ 坐标系的第一象限内，曲线 $\eta = -\xi \cot(\xi)$ 与圆 $\xi^2 + \eta^2 = R^2$ 至少存在一个交点的条件是什么？}

我们分析曲线 $\eta = -\xi \cot(\xi)$ 的性质：
\begin{itemize}
    \item 当 $0 < \xi < \frac{\pi}{2}$ 时，$\cot(\xi) > 0$，故 $\eta = -\xi \cot(\xi) < 0$。这说明曲线在第一象限（$\xi \in (0, \pi/2)$）内\textbf{无图象}（因为要求 $\eta > 0$）。
    \item 当 $\frac{\pi}{2} < \xi < \pi$ 时，$\cot(\xi) < 0$，故 $\eta = -\xi \cot(\xi) > 0$。曲线在此区间进入上半平面（即第一象限所在区域）。
    \item 当 $\xi \to \left(\frac{\pi}{2}\right)^+$ 时，$\cot(\xi) \to 0^-$，从而 $\eta \to 0^+$。这表明曲线与 $\xi$ 轴的正半轴交于点 $\left(\frac{\pi}{2}, 0\right)$，并随着 $\xi$ 增大向右上方单调递增。
\end{itemize}

因此，为了让半径为 $R$ 的圆心在原点的圆弧，能够与从 $\left(\frac{\pi}{2}, 0\right)$ 出发的曲线 $\eta = -\xi \cot(\xi)$ 在第一象限（严格来说是 $\eta > 0$ 区域）发生相交，圆的半径 $R$ 必须严格大于交点的横坐标起点：
\[
    R > \frac{\pi}{2}
\]
将 $R$ 的表达式代回：
\[
    \sqrt{\frac{2mV_0a^2}{\hbar^2}} > \frac{\pi}{2}
\]
两边平方并整理，得到势阱深度 $V_0$ 和宽度 $a$ 必须满足的关系：
\[
    V_0 a^2 > \frac{\pi^2 \hbar^2}{8m}
\]

\end{solution}
\end{document}